\documentclass{article}


% if you need to pass options to natbib, use, e.g.:
%     \PassOptionsToPackage{numbers, compress}{natbib}
% before loading neurips_2023

\PassOptionsToPackage{authoryear,round}{natbib}

% ready for submission
\usepackage[final]{neurips_2023}


% to compile a preprint version, e.g., for submission to arXiv, add add the
% [preprint] option:
%     \usepackage[preprint]{neurips_2023}


% to compile a camera-ready version, add the [final] option, e.g.:
%     \usepackage[final]{neurips_2023}


% to avoid loading the natbib package, add option nonatbib:
%    \usepackage[nonatbib]{neurips_2023}


\usepackage[utf8]{inputenc} % allow utf-8 input
\usepackage[T1]{fontenc}    % use 8-bit T1 fonts
\usepackage{hyperref}       % hyperlinks
\usepackage{url}            % simple URL typesetting
\usepackage{booktabs}       % professional-quality tables
\usepackage{amsfonts}       % blackboard math symbols
\usepackage{nicefrac}       % compact symbols for 1/2, etc.
\usepackage{microtype}      % microtypography
\usepackage{xcolor}         % colors

\usepackage{graphicx}
\usepackage{caption}
\usepackage{float}
\captionsetup[figure]{font=footnotesize}

\title{Geometric Deep Learning}


% The study will b e assessed based o n the following criteria: (i) motivation, clarity, and presenta-
% tion (20 marks), (ii) originality and novelty (20 marks), (iii) coherence and depth o f the study
% (20 marks), (iv) scholarship, i.e., framing i n the existing literature (15 marks) (v) relations to
% the concepts discussed throughout the course (10 marks), (vi) balanced critical self-evaluation
% (15 marks).
% The study will not be evaluated according t o the length, number of experiments or datasets, but
% according t o its merits in accordance to the outlined criteria. The following will be the rough
% grading guidance: 90-100: exceptionally written papers worth publishing that show novel theo-
% retical results, propose a novel method or a significant improvement of an existing one, or report
% state-of-the-art experimental results; 80-89: very well-written papers that show some theoreti-
% cal novelty, incremental improvement of a method, or a small improvement of state-of-the-art
% experimental results; 70-79: well-written papers that show a good reproduction of an existing
% method but lack insights on drawbacks and possible theoretical or experimental improvements.

% The report should not exceed nine pages and should adhere to NeurIPS style
% (https://neurips.cc/Conferences/2023/PaperInformation/StyleFiles). Additional pages contain-
% ing only the references are allowed. Appendices (showing e.g. derivations and implementation
% details) are allowed but should not b e necessary for the reading and evaluation of the report.
% Please follow the notation from the lecture slides whenever applicable. If the study involves
% a n implementation, please provide your code with the submission as a zip file. 

\author{Candidate Number 1091797}

\begin{document}

\maketitle

\begin{abstract}
    Hello
\end{abstract}

\section{Introduction}
% State a research question related t o geometric deep learning and explain your motivation for the proposed study. (Example: How can the use of graph direction help message passing?)
\subsection{Problem statement}
\subsection{Motivation}
\subsection{Background}

\section{Outline of our study and goals}

% Propose means to study the research question and give an outline of the overall approach,
% clearly identifying the goals of the chosen study. (Example: An experimental setup aiming
% to test different ways o f directed message passing o n graphs, using a synthetic dataset or
% using existing datasets.)

\section{Methodology}

% Clearly state your methodology: theoretical framework and empirical setup, hyper-
% parameters, and the assumptions underpinning the study.

\section{Experiments}

% Report your empirical/theoretical/conceptual findings, identifying whether o r not the re-
% sults support the initial hypothesis. You can use visuals, figures, and tables to present
% your findings in a structured manner. You are expected to relate and compare your results
% to the relevant literature and t o the concepts from the course.

\section{Discussion}

% Provide a detailed discussion relating the results t o the original motivation, as well as a
% critical perspective regarding the study.

\section{Outlook} 

% The report should conclude with an outlook stating any additional studies that need to be conducted to reach more conclusive statements


% \newpage
{
% \small 
% \bibliographystyle{plainnat}
\bibliography{references}
}

\end{document}